\section[Application to principal coverings]{Application to principal coverings: Kummer and
Artin-Schreier theories}
\label{XI.6}

\begin{proposition}
\label{XI.6.1}
Let $S$ be a prescheme, $n$ an integer $>0$, let $u_n\colon
\GG_{m\,S}\to\GG_{m\,S}$ be the $n$th power homomorphism, and
$\bbmu_{n\,S}$ its kernel. Then $\bbmu_{n\,S}$ is finite and locally
free of rank $n$ over~$S$, and it is \'etale over~$S$ if and only if
for every $s\in S$, the characteristic of~$s$ is prime to
$n$. The sequence of homomorphisms
$$
0\to \bbmu_{n\,S}\to\GG_{m\,S}\lto{u_n}\GG_{m\,S}\to 0
$$
is exact in the sense of \No \Ref{XI.4}. (It will be called the \emph{Kummer
exact sequence} over~$S$, relative to the integer~$n$).
\end{proposition}

We
% original p. 303
have
$$
\GG_m=\Spec\cal{O}_S[t,t^{-1}],
$$
and $u_n$ corresponds to the homomorphism $u_n$ on the affine
$\cal{O}_S$-algebras, given by
$$
u_n(t)=t^n,
$$
on the other hand the unit section of~$\GG_{m\,S}$ corresponds to
the augmentation homomorphism of~$\cal{O}_S$-algebras, given
by
$$
\varepsilon(t)=1,
$$
whose kernel is therefore the principal Ideal $(t-1)$. The image of the
latter under $u_n$ is therefore the principal Ideal $(1-t^n)$, and one
finds:
$$
\bbmu_{n\,S}=\Spec\cal{O}_S[t]/(1-t^n),
$$
which shows in particular that $\bbmu_{n\,S}$ is finite over~$S$,
and defined by an Algebra over~$S$ which is free of rank $n$,
having the basis formed by the $t^i$ ($0\leq i\leq n-1$). For it to be
\'etale at $s\in S$, it is necessary and sufficient that the reduced
Algebra $k[t]/(1-t^n)$, where $k=\kres(s)$, obtained by formal
adjunction of the $n$th roots of unity to $k$, be
separable over~$k$, \ie the roots of~$1-t^n$ in an algebraic
closure of~$k$ are all distinct, which is equivalent to the fact
that $n$ is prime to the characteristic. Finally, to show that the
sequence of homomorphisms in~\Ref{XI.6.1} is exact, one is reduced,
by virtue of criterion~\Ref{XI.4.2}, to proving that
$u_n$
is faithfully flat. One can obviously assume $S$ affine
with ring~$A$, hence $\GG_{m\,S}$ affine with ring $B=A[t,t^{-1}]$, and it
suffices to verify that $u_n$ makes $B$ a free module of rank
$n$ over~$B$, or, what amounts to the same, that $u_n$ is injective, and
that $A[t,t^{-1}]$ is a free module of rank $n$ over~$A[t^n,t^{-n}]$.
Indeed, one easily verifies that the
$t^i$ ($0\leq i\leq n-1$) form a basis of the one over the other, which
completes the proof.

\begin{definition}
\label{XI.6.2}
We
% original p. 304
call $\bbmu_{n\,S}$ the \emph{Kummer group of rank $n$
over~$S$}, and we call \emph{Kummer principal covering
of rank
$n$ over~$S$}
every principal homogeneous bundle over~$S$ with group the
Kummer group of rank~$n$.
\end{definition}

The set of these coverings is a group, denoted
$\H^1(S,\bbmu_{n\,S})$ or simply $\H^1(S,\bbmu_n)$. Note that
the formation of the Kummer group of rank $n$ over~$S$ is
compatible with base extension, hence that $\bbmu_{n\,S}$
comes by base extension from the \emph{absolute Kummer group
$\bbmu_n$} over~$\Spec(\ZZ)$.

Let us denote by $(\ZZ/n\ZZ)_S$ the $S$-group defined by the
ordinary finite group $\ZZ/n\ZZ$. If~$G$ is an arbitrary $S$-group,
the homomorphisms of~$S$-groups $u$ from~$(\ZZ/n\ZZ)_S$ to $G$
correspond one-to-one, and compatibly with
change of base, to the sections of~$G$ over~$S$ whose $n$th power
is the unit section, by associating with
$u$ the image under $u$ of the section of~$(\ZZ/n\ZZ)_S$ over~$S$ defined
by the generator $1\bmod n\ZZ$ of~$\ZZ/n\ZZ$. This being said:

\begin{corollary}
\label{XI.6.3} If $\bbmu_{n\,S}$ is \'etale over~$S$, one
thus obtains a one-to-one correspondence between the isomorphisms of~$S$-groups
$(\ZZ/n\ZZ)_S\isomto\bbmu_{n\,S}$, and the sections of~$\cal{O}_S$
which are of order exactly $n$ on each connected component
of~$S$ (such a section will be called a ``primitive
$n$th root of unity over~$S$''). Hence for
$\bbmu_{n\,S}$ to be isomorphic as an $S$-group to
$(\ZZ/n\ZZ)_S$, it is necessary and sufficient that it be \'etale
over~$S$,
\ie that the residue characteristics of~$S$ be
prime to $n$, and that there exist a primitive
$n$th root of unity over~$S$.
\end{corollary}

This explains the role played in classical Kummer theory
by the hypothesis that the base field (playing the role
of~$S$) be of characteristic prime to $n$ and
contain the $n$th roots of unity, and by the choice
of a primitive $n$th root of unity. Once one
has the language of schemes at one's disposal, there is no longer any reason to
encumber oneself with these hypotheses, and one should reason directly
on~$\bbmu_n$ instead of~$\ZZ/n\ZZ$. Thus, the conjunction
of~\Ref{XI.6.1}, \Ref{XI.4.5} and~\Ref{XI.5.3} gives us the following
general relation between the theory of Kummer principal
coverings and that of Picard groups:

\begin{proposition}
\label{XI.6.4}
Let
% original p. 305
$S$ be a prescheme; one has a canonical exact sequence
$$
\textstyle
0\to\H^0(S,\bbmu_n)\to\H^0(S,\cal{O}^*_S)\to
\H^0(S,\cal{O}^*_S)\to\H^1(S,\bbmu_n)\to
\H^1(S,\cal{O}^*_S)\to\H^1(S,\cal{O}^*_S),
$$
hence, setting
$\H^1(S,\cal{O}_S^*)=\Pic(S)$,
and denoting, for every abelian group $A$, by ${_n}A$ and
$A_n$ the kernel and cokernel of multiplication by $n$ in $A$, the
exact sequence:
$$
0\to\H^0(S,\cal{O}_S)^*_n\to\H^1(S,\bbmu_n)\to
{\sideset{_n}{}\Pic}(S)\to 0.
$$
\end{proposition}

We are going to make explicit two important cases, where one or the other
extreme term of this exact sequence is zero:

\begin{corollary}
\label{XI.6.5} Suppose $\sideset{_n}{}\Pic(S)=0$, (for example
that $S$ is the spectrum of a local ring, or of a factorial ring),
and let $A$ be the ring $\H^0(S,\cal{O}_S)$. Then one has a canonical
isomorphism
$$
\H^1(S,\bbmu_n)\isomto A^*/{A^*}^n.
$$
\end{corollary}

This is essentially the classical statement of Kummer theory,
when $S$ is the spectrum of a field.

\begin{corollary}
\label{XI.6.6} Suppose that every element of
$\H^0(S,\cal{O}_S)$ is an $n$th power, for example that
$\H^0(S,\cal{O}_S)$ is a composite of algebraically
closed fields or that $S$ is reduced and proper over an
algebraically closed field $k$. Then one has a canonical isomorphism
$$
\H^1(S,\bbmu_n)\isomto\sideset{_n}{}\Pic(S).
$$
\end{corollary}

In particular, when $S$ is proper and connected over an
algebraically closed field $k$, this relates the fundamental
group of~$S$ to the points of finite order of the Picard
scheme $P$ of~$S$ over~$k$; thus one will have an isomorphism
$$
\Hom(\pi_1(S),\ZZ/n\ZZ)\simeq {_n P}(k)
$$
for $n$ prime to the characteristic, which is often
used in algebraic
% original p. 306
geometry. As an application, when the connected component $P^0$
of~$P$ is a complete group scheme, of dimension $g$, one sees,
using the results recalled in \No \Ref{XI.2} and the
finiteness of the torsion group of
N\'eron-Severi,
that for
every prime number $\ell$ prime to the characteristic, the
$\ell$-primary component of the fundamental group $\pi_1(S)$ made
abelian is a module of finite type and of rank $2g$ over the ring
$\ZZ_\ell$ of $\ell$-adic integers (and moreover free except for
at most a finite number of values of~$\ell$). As
Serre has remarked, this makes it possible to prove, under certain conditions, that when
$X$ is a flat and projective scheme over a connected~$S$, then the
Picard schemes of the fibers of~$X$ all have the same
dimension, by applying the semicontinuity theorem
(X~\Ref{X.2.3}); Serre's argument applies as soon as the
Picard scheme of~$X$ over~$S$ exists and the connected
Picards of the fibers of~$X$ over~$S$ are proper group
schemes, for example when the geometric fibers of~$X$
over~$S$ are normal ($X$ still being flat and projective
over~$S$), in particular if $X$ is smooth and projective over~$S$.

Now let $p$ be a prime number, and suppose that $S$ is a
prescheme of characteristic $p$, \ie such that
$p\cdot\cal{O}_S=0$. Then the $p$th power homomorphism
in $\cal{O}_S$ is additive, and the corresponding morphism, obtained by
replacing $S$ by a variable~$T$ over~$S$:
$$
F\colon\GG_{a\,S}\to\GG_{a\,S}
$$
is therefore a homomorphism of~$S$-groups, which is called
the \emph{Frobenius homomorphism} (N.B. Such a morphism is
defined for every $S$-prescheme $G$ which comes by
base extension from a prescheme $G_0$ over the prime field
$\ZZ/p\ZZ$, and this morphism is a group homomorphism if $G_0$
is a group prescheme). We shall set:
$$
\wp=\id-F\colon\GG_{a\,S}\to\GG_{a\,S}.
$$
Consider on the other hand the $S$-group $(\ZZ/p\ZZ)_S$ defined
by the ordinary finite group~$\ZZ/p\ZZ$; we have said that for every
$S$-group~$G$, the homomorphisms of~$S$-groups from~$(\ZZ/p\ZZ)_S$
to $G$ correspond one-to-one to the sections of~$G$ over~$S$ whose
$p$th power is the unit section. When
$G=\GG_{a\,S}$, they
% original p. 307
therefore correspond to arbitrary sections of~$G$ over~$S$. Taking in particular
the section of~$\GG_{a\,S}$ over~$S$ corresponding to the unit
section of the sheaf of rings $\cal{O}_S$, one finds a
homomorphism of~$S$-groups
$$
i\colon(\ZZ/p\ZZ)_S\to\GG_{a\,S}
$$

\begin{proposition}
\label{XI.6.7} The sequence of homomorphisms of~$S$-groups
$$
0\to(\ZZ/p\ZZ)_S\to\GG_{a\,S}\to \GG_{a\,S}\to 0
$$
is exact (in the sense of \No \Ref{XI.4}). (It is called the
\emph{Artin-Schreier} exact sequence over~$S$).
\end{proposition}

It suffices to prove it over the prime field $k=\ZZ/p\ZZ$. It suffices
to note that the homomorphism $\wp^*\colon k[t]\to k[t]$ defined
by $\wp^*(t)=t-t^p$ makes $k[t]$ a free module of rank $p$ over~$k[t]$,
more precisely that $k[t]$ is a free module over~$k[s]$,
where $s=t-t^p$, having the basis formed by the $t^i$ ($0\leq i\leq p-1$).

One concludes from this, using \Ref{XI.4.5} and~\Ref{XI.5.3}:

\begin{proposition}
\label{XI.6.8} One has a canonical exact sequence:
\begin{multline*}
0\to\H^0(S,\ZZ/p\ZZ)\to\H^0(S,\cal{O}_S)\to\H^0(S,\cal{O}_S)\to
\H^1(S,\ZZ/p\ZZ)\\
\to\H^1(S,\cal{O}_S)\to\H^1(S,\cal{O}_S),
\end{multline*}
hence an exact sequence:
$$
0\to\H^0(S,\cal{O}_S)/\wp\H^0(S,\cal{O}_S)\to
\H^1(S,\ZZ/p\ZZ)\to\H^1(S,\cal{O}_S)^F\to 0,
$$
(where the exponent $F$ in the last term means the subgroup
of invariants under the endomorphism $F$, equal to the kernel of
$\wp=\id-F$).
\end{proposition}

Let us make explicit two more extreme cases:

\begin{corollary}
\label{XI.6.9} Suppose that $\H^1(S,\cal{O}_S)^F=0$, for example
that $S$ is an affine scheme. Then, setting
$A=\H^0(S,\cal{O}_S)$, one has a canonical isomorphism
$$
\H^1(S,\ZZ/p\ZZ)\isomto A/\wp A.
$$
\end{corollary}

This is \emph{Artin-Schreier theory} in the classical form,
at least when $A$ is the spectrum of a field.

\begin{corollary}
\label{XI.6.10}
Suppose
% original p. 308
that $\wp\H^0(S,\cal{O}_S)\mkern1mu{=}\mkern1mu\H^0(S,\cal{O}_S)$, for example
that $\H^0(S,\cal{O}_S)$ is a composite of algebraically closed
fields, or that $S$ is proper over an algebraically closed
field. Then one has a canonical isomorphism:
$$
\H^1(S,\ZZ/p\ZZ)\isomto\H^1(S,\cal{O}_S)^F
$$
\end{corollary}

\begin{remarks}
\label{XI.6.11}
The last statement is due to
J-P\ptbl Serre
\cite{XI.9}. It is
also possible to develop an analogous theory for
the structure group $\ZZ/p^n\ZZ$ for arbitrary $n$, using
instead of~$\GG_a$ the Witt group scheme $\WW_n$, \cf \loccit Note
that in characteristic $p>0$, Kummer theory no longer gives
any information on principal coverings of order $p$, since
$\bbmu_p$ is then an ``infinitesimal'' group \ie radicial over the
base, hence without direct relation to $\ZZ/p\ZZ$; so at first sight,
the theory of these coverings is no longer amenable (when $S$ is a
proper scheme over an algebraically closed field, to fix
ideas) to the theory of the Picard scheme as
in~\Ref{XI.6.6}. Nevertheless, if one recalls that the Zariski tangent
space at the origin in
$\mathbf{Pic}_{S/K}$\kern1pt\footnote{For the definition of
$\mathbf{Pic}_{S/K}$, \cf A. Grothendieck, S\'em. Bourbaki \No 232
(February 1962).} is identified with $\H^1(S,\cal{O}_S)$, one
observes that \emph{knowledge of the group scheme
${_p}\mathbf{Pic}_{S/k}$, kernel of multiplication by $p$ in
$\mathbf{Pic}_{S/k}$, implies that of~$\H^1(S,\ZZ/p\ZZ)$ as well
as that of~$\H^1(S,\bbmu_p)$; note that it also implies
that of~$\H^1(S,\bbalpha_p)$}, where $\bbalpha_p$
denotes the infinitesimal group scheme over the prime
field, kernel of~$F\colon\GG_a\to\GG_a$ (which can also be described
as the spectrum of the restricted enveloping algebra of the
trivial $p$-Lie algebra of dimension~$1$): indeed the exact
sequence~\Ref{XI.4.5} gives here:
$$
\H^1(S,\bbalpha_p)\simeq\Ker
(F\colon\H^1(S,\cal{O}_S)\to\H^1(S,\cal{O}_S)),
$$
and more generally, denoting by $\bbalpha_{p^n}$
the kernel in $\GG_a$ of the $n$th iterate of~$F$, one will have
$$
\H^1(S,\bbalpha_{p^n})\simeq\Ker(F^n\colon\H^1(S,\cal{O}_S)\to
\H^1(S,\cal{O}_S)).
$$

In fact, knowledge of~${_p}\mathbf{Pic}_{S/k}$ is equivalent to
that of~$\H^1(S,G)$ for every finite commutative algebraic group
% original p. 309
annihilated by $p$; more generally,
knowledge of~${_{p^n}}\mathbf{Pic}_{S/k}$ is equivalent to that
of~$\H^1(S,G)$ for every finite commutative algebraic group $G$
annihilated by $p^n$, by virtue of the following theorem, which encompasses,
in the case under consideration, both Kummer theory and
that of Artin-Schreier:

Let $G$ be a finite algebraic group over~$k$, and
$D(G)=\SheafHom_{k\textup{-groups}}(G,\GG_m)$ its \emph{Cartier
dual} (whose affine algebra is carried by the vector space
dual to the affine algebra of~$G$, \ie by
the hyperalgebra of~$G$ in the sense of Dieudonn\'e-Cartier); then one
has a canonical isomorphism:
\begin{equation}\label{eqXI.6.11}\tag{$*$}
\H^1(S,G)\simeq\Hom_{k\textup{-groups}}(D(G),\mathbf{Pic}_{S/k}).
\end{equation}
(N.B. $S$ is a proper scheme over~$k$ algebraically closed, such
that $\H^0(S,\cal{O}_S)=k$). This formula can also be expressed by
saying that the ``true fundamental group'' of~$S$ alluded to
in \No \Ref{XI.2}, made abelian, is isomorphic to the
projective limit of the $D(P_i)$, where $P_i$ runs through the
\emph{finite} algebraic subgroups of~$\mathbf{Pic}_{S/k}$, which will be denoted
$T^\bbullet(\mathbf{Pic}_{X/k})$. When $S$ is an abelian
variety, we have seen in~\Ref{XI.2.1} that this group is
also isomorphic to the ``true'' Tate module
$T_{\bbullet}(S)=\varprojlim\, {{_n}S}$, and the isomorphism
\eqref{eqXI.6.11} can then be written more strikingly as
$$
\Ext^1(A,G)\simeq\Hom(D(G),B),
$$
$A$ being an abelian variety, $B$ its dual, $G$ a
finite algebraic group over~$k$. The results just
indicated can moreover be generalized to the case where $k$
is replaced by an arbitrary base prescheme, and to
coefficient groups $G$ other than finite groups.
\end{remarks}

\begin{thebibliography}{10}{XI.7}
% original p. 310

\bibitem[1]{XI.1} A\ptbl \textsc{Grothendieck},
Sur quelques points d'alg\`ebre homologique,
T\^ohoku Math.~J.\ \textbf{9} (1957) p\ptbl 119--221.

\bibitem[2]{XI.2} A\ptbl \textsc{Grothendieck}, A general theory of fibre spaces with
structure sheaf, University of Kansas, 1955.

\bibitem[3]{XI.3} A\ptbl \textsc{Grothendieck}, Torsion homologique et sections
rationnelles, S\'eminaire \hbox{Chevalley}, June 16, 1958.

\bibitem[4]{XI.4} S\ptbl \textsc{Lang}, Abelian varieties, Interscience tracts in pure
and applied mathematics, \No 7, New York.

\bibitem[5]{XI.5}
J-P\ptbl%
\textsc{Serre}, Groupes alg\'ebriques et corps de classes,
Actualit\'es Scient. et Ind.\ \No 1264, Hermann, Paris, 1959.

\bibitem[6]{XI.6}
J-P\ptbl%
\textsc{Serre}, Groupes proalg\'ebriques,
Publications Math\'ematiques de l'IHES \textbf{7} (1960), p\ptbl 1--67.

\bibitem[7]{XI.7}
J-P\ptbl%
\textsc{Serre}, Espaces fibr\'es alg\'ebriques,
S\'eminaire Chevalley, April 21, 1958.

\bibitem[8]{XI.8}
J-P\ptbl%
\textsc{Serre}, Quelques propri\'et\'es des
vari\'et\'es ab\'eliennes en
caract\'eristique~$p$, Amer.\ J.\ Math. \textbf{80} (1958), p\ptbl 715--739.

\bibitem[9]{XI.9}
J-P\ptbl%
\textsc{Serre}, Sur la topologie des vari\'et\'es
alg\'ebriques en
caract\'eristique~$p$,
Symposium Internacional de Topologia
Algebraica, 1958.

\bibitem[10]{XI.10}
J-P\ptbl%
\textsc{Serre}, On the fundamental group of a unirational
variety,
J.~London Math.\ Soc.\ \textbf{34} (1959), p\ptbl 481--484.
\end{thebibliography}
